% !TeX root = anexo_semana03.tex
% !TeX program = pdflatex
% Material complementario de la semana 03: funciones seno y coseno.
\documentclass[10pt,a5paper]{article}
\newif\ifmostrarSoluciones
\mostrarSolucionestrue
\usepackage{tikz}
\usepackage{estilo_apuntes}
\usetikzlibrary{arrows.meta,positioning,babel}
\encabezado{Semana 3}{Material complementario}

\begin{document}
\tituloclase{Las funciones seno y coseno}{Definición geométrica, dominio, recorrido y gráficas}{Material complementario · Semana 3}

\begin{objetivos}
\begin{itemize}
  \item Interpretar el seno y el coseno como coordenadas de un punto del círculo unitario.
  \item Identificar su dominio, recorrido, periodicidad y cotas.
  \item Reconocer estas propiedades en las gráficas de ambas funciones.
\end{itemize}
\end{objetivos}

\begin{resumen}
Las funciones seno y coseno permiten describir un movimiento alrededor del círculo unitario. Sus gráficas muestran que ambas funciones son periódicas y que sus valores siempre permanecen entre $-1$ y $1$.
\end{resumen}

\section*{Definición mediante el círculo unitario}

Los ángulos se medirán en \textbf{radianes}. Para cada $t\in\mathbb{R}$, partimos del punto $(1,0)$ y recorremos una distancia $|t|$ sobre el círculo unitario, en sentido antihorario si $t>0$ y horario si $t<0$. Si el punto obtenido es
\[
  P(t)=(u,v),
\]
definimos
\[
  \cos(t)=u
  \qquad\text{y}\qquad
  \sin(t)=v.
\]
Así, el coseno es la primera coordenada de $P(t)$ y el seno es la segunda.

La siguiente figura representa el caso $0<t<\pi/2$. Las líneas punteadas proyectan el punto $P(t)$ sobre los ejes y permiten leer sus dos coordenadas.
\begin{center}
\begin{tikzpicture}[x=2.15cm,y=2.15cm,every node/.style={font=\scriptsize}]
  \draw[->,Gris] (-1.25,0)--(1.35,0) node[right] {$u$};
  \draw[->,Gris] (0,-1.25)--(0,1.35) node[above] {$v$};
  \draw[Azul!75,thick] (0,0) circle (1);
  \node[Azul!75,below left] at (-0.72,-0.70) {$u^2+v^2=1$};
  \coordinate (P) at (50:1);
  \coordinate (A) at (0.6428,0);
  \coordinate (B) at (0,0.7660);
  \draw[Rojo,very thick] (0,0)--(P);
  \draw[Gris,dashed] (P)--(A);
  \draw[Gris,dashed] (P)--(B);
  \draw[->,Rojo,thick] (0.30,0) arc[start angle=0,end angle=50,radius=0.30];
  \node[Rojo] at (0.38,0.16) {$t$};
  \fill[Rojo] (P) circle (1.5pt);
  \node[above right] at (P) {$P(t)=(\cos(t),\sin(t))$};
  \fill[Gris] (1,0) circle (1.1pt);
  \node[below right] at (1,0) {$(1,0)$};
  \node[below] at (A) {$\cos(t)$};
  \node[left] at (B) {$\sin(t)$};
  \node[below left] at (0,0) {$O$};
\end{tikzpicture}
\end{center}

\begin{definicion}[seno y coseno]
Las funciones seno y coseno son
\[
  \sin\colon\mathbb{R}\longrightarrow[-1,1],
  \qquad
  \cos\colon\mathbb{R}\longrightarrow[-1,1].
\]
Ambas tienen dominio $\mathbb{R}$ y recorrido $[-1,1]$.
\end{definicion}

Al dar una vuelta completa se recorren $2\pi$ radianes y se regresa al mismo punto. Por tanto,
\[
  \sin(t+2\pi)=\sin(t),
  \qquad
  \cos(t+2\pi)=\cos(t).
\]
Además, como las coordenadas de un punto del círculo unitario están entre $-1$ y $1$,
\[
  -1\leq\sin(t)\leq1,
  \qquad
  -1\leq\cos(t)\leq1
  \qquad (t\in\mathbb{R}).
\]
Los valores correspondientes a una vuelta completa son:
\[
\begin{array}{c|ccccc}
  t & 0 & \dfrac{\pi}{2} & \pi & \dfrac{3\pi}{2} & 2\pi \\ \hline
  \sin(t) & 0 & 1 & 0 & -1 & 0 \\
  \cos(t) & 1 & 0 & -1 & 0 & 1
\end{array}
\]

\section*{Gráficas}

\begin{center}
\begin{tikzpicture}[x=0.68cm,y=1.05cm,every node/.style={font=\scriptsize}]
  \draw[->,Gris] (-6.7,0)--(6.8,0) node[right] {$t$};
  \draw[->,Gris] (0,-1.35)--(0,1.4) node[above] {$y$};
  \draw[Gris!55,dashed] (-6.4,1)--(6.4,1);
  \draw[Gris!55,dashed] (-6.4,-1)--(6.4,-1);
  \draw[Azul,very thick,domain=-6.283:6.283,samples=220]
    plot (\x,{sin(deg(\x))});
  \foreach \x/\etiqueta in {-6.283/$-2\pi$,-3.142/$-\pi$,3.142/$\pi$,6.283/$2\pi$}
    \draw (\x,0.06)--(\x,-0.06) node[below] {\etiqueta};
  \node[left] at (0,1) {$1$};
  \node[left] at (0,-1) {$-1$};
  \node[Azul,above right] at (4.7,-1) {$y=\sin(t)$};
\end{tikzpicture}
\end{center}

\begin{center}
\begin{tikzpicture}[x=0.68cm,y=1.05cm,every node/.style={font=\scriptsize}]
  \draw[->,Gris] (-6.7,0)--(6.8,0) node[right] {$t$};
  \draw[->,Gris] (0,-1.35)--(0,1.4) node[above] {$y$};
  \draw[Gris!55,dashed] (-6.4,1)--(6.4,1);
  \draw[Gris!55,dashed] (-6.4,-1)--(6.4,-1);
  \draw[Rojo,very thick,domain=-6.283:6.283,samples=220]
    plot (\x,{cos(deg(\x))});
  \foreach \x/\etiqueta in {-6.283/$-2\pi$,-3.142/$-\pi$,3.142/$\pi$,6.283/$2\pi$}
    \draw (\x,0.06)--(\x,-0.06) node[below] {\etiqueta};
  \node[left] at (0,1) {$1$};
  \node[left] at (0,-1) {$-1$};
  \node[Rojo,above right] at (5.0,0.25) {$y=\cos(t)$};
\end{tikzpicture}
\end{center}

Las gráficas muestran que ambas funciones permanecen entre $-1$ y $1$, y que sus valores se repiten indefinidamente.

\section*{La función $\sin(1/x)$}

Considere la función
\[
  f\colon\mathbb{R}\setminus\{0\}\longrightarrow[-1,1],
  \qquad
  f(x)=\sin(1/x).
\]
La función no está definida en $x=0$. Sin embargo, para todo punto $x$ en el dominio de la función se cumple
\[
  -1\leq\sin(1/x)\leq1.
\]
Por tanto, la función permanece acotada entre $-1$ y $1$.

\begin{center}
\begin{tikzpicture}[x=4.1cm,y=1.55cm,every node/.style={font=\scriptsize}]
  \draw[->,Gris] (-1.08,0)--(1.10,0) node[right] {$x$};
  \draw[->,Gris] (0,-1.35)--(0,1.38) node[above] {$y$};
  \draw[Rojo,dashed] (-1.03,1)--(1.03,1);
  \draw[Rojo,dashed] (-1.03,-1)--(1.03,-1);
  \draw[Azul,thick,domain=-1:-0.015,samples=900]
    plot (\x,{sin(deg(1/\x))});
  \draw[Azul,thick,domain=0.015:1,samples=900]
    plot (\x,{sin(deg(1/\x))});
  \foreach \x/\etiqueta in {-1/$-1$,-0.5/$-0.5$,0.5/$0.5$,1/$1$}
    \draw (\x,0.04)--(\x,-0.04) node[below] {\etiqueta};
  \node[left,Rojo] at (0,1) {$1$};
  \node[left,Rojo] at (0,-1) {$-1$};
  \node[below right] at (0,-0.08) {$x=0$};
  \node[Azul,above right] at (0.62,1) {$y=\sin(1/x)$};
\end{tikzpicture}
\end{center}

La altura de la gráfica nunca sale de la franja comprendida entre $y=-1$ y $y=1$. En cambio, las oscilaciones se hacen cada vez más rápidas al acercarse a $x=0$: un pequeño cambio en $x$ produce un cambio grande en $1/x$ y hace que el seno recorra repetidamente sus valores entre $-1$ y $1$. Estar acotada no significa que la función se aproxime a un único número.

\fuentes{J. Stewart, \emph{Cálculo de una variable: trascendentes tempranas}, apartados sobre funciones trigonométricas y límites.}

\end{document}
